\documentclass[11pt]{article}

\usepackage[a4paper,margin=0.82in]{geometry}
\usepackage[T1]{fontenc}
\usepackage[utf8]{inputenc}
\usepackage{lmodern}
\usepackage{amsmath,amssymb,amsthm,mathtools}
\usepackage{booktabs,enumitem,microtype,tabularx,xcolor}
\usepackage[numbers,sort&compress]{natbib}
\usepackage[colorlinks=true,linkcolor=blue!65!black,
 citecolor=blue!65!black,urlcolor=blue!70!black]{hyperref}
\usepackage[nameinlink,noabbrev]{cleveref}
\hypersetup{
 pdftitle={Complete High-Frequency and Ultra-Tail Closure at the Original Physical Normalization},
 pdfauthor={Liang Wang},
 pdfsubject={TPC-116 exact tail partition, sharp Hilbert criterion, and phase-blind obstruction},
 pdfkeywords={high-frequency tail, physical normalization, exact partition, phase obstruction, fixed shift}
}

\newtheorem{theorem}{Theorem}[section]
\newtheorem{proposition}[theorem]{Proposition}
\newtheorem{corollary}[theorem]{Corollary}
\newtheorem{lemma}[theorem]{Lemma}
\theoremstyle{definition}
\newtheorem{definition}[theorem]{Definition}
\theoremstyle{remark}
\newtheorem{remark}[theorem]{Remark}

\newcommand{\C}{\mathbb C}
\newcommand{\one}{\mathbf 1}
\newcommand{\Lzero}{\mathrm{L0}}
\newcommand{\Lone}{\mathrm{L1}}
\newcommand{\Ltwo}{\mathrm{L2}}
\newcommand{\Soft}{\mathsf S}
\newcommand{\Ret}{\mathsf R}
\newcommand{\Tail}{\mathcal T}
\newcommand{\Low}{\mathrm{Low}}
\newcommand{\High}{\mathrm{High}}
\newcommand{\Ultra}{\mathrm{Ultra}}
\newcommand{\Short}{\mathrm{Short}}
\newcommand{\Boundary}{\mathrm{Boundary}}
\newcommand{\Collision}{\mathrm{Collision}}

\title{\textbf{Complete High-Frequency and Ultra-Tail Closure}\\
\textbf{at the Original Physical Normalization:}\\
An Exact Partition, an Optimal Hilbert Test,\\
and a Sharp Phase-Blind Obstruction}
\author{Liang Wang\\
School of Artificial Intelligence and Automation\\
Huazhong University of Science and Technology, Wuhan 430074, P.R. China\\
\texttt{wangliang.f@gmail.com}}
\date{July 2026}

\begin{document}
\maketitle

\begin{abstract}
The TPC reconstruction program has several rigorous low-window
interfaces, but its working Hypothesis H4 asks for more: every
high-frequency, ultra-long, short-fiber, boundary, collision and
polarization complement must return at the \emph{original physical
normalization}.  A tail estimate in one row or one favorable sector
does not by itself meet that quantifier.

We give a finite exact closure theorem conditional on a declared
complete mask family.  We characterize exact one-fold partition of
the literal physical atom set into the low window and all complement
classes.  On a frequency class, unitary Fourier
coordinates give the optimal estimate
\[
 |\Tail_E|
 \le \nu_X
 \Big(\sum_{r\in E}|\widehat m_X(r)|^2\Big)^{1/2}
 \Big(\sum_{r\in E}|\widehat A_X(r)|^2\Big)^{1/2};
\]
the constant one is sharp.  Combining this with complete outer
absolute-mass bounds for soft classes and literal aggregate packet
bounds for retained classes yields a sufficient original-scale
closure.  Every reconstruction multiplier, provenance multiplicity
and outer normalization must occur inside those returned bounds.

We also prove a sharp obstruction.  If only the magnitudes of the
physical atoms are known, the worst coherent choice of phases has
size exactly their total absolute mass.  Hence support sparsity,
shortness of individual fibers, or a partition into many blocks
cannot yield a power saving without coefficient-sensitive
cancellation or an absolute-mass theorem.  Existing TPC low-window
tail reserves therefore remain valid local inputs, but H4 is not
closed until their complete outer transport and every ultra-tail
class are certified.  The result is exact \(\Lzero/\Lone\)
bookkeeping; it proves no new \(\Ltwo\) fixed-shift estimate, parity
breakthrough or twin-prime theorem.
\end{abstract}

\noindent\textbf{Keywords:}
high-frequency complement; ultra-tail; physical normalization;
Fourier energy; phase alignment; stop--go certificate.

\medskip
\noindent\textbf{Physical convention.}
One fixed nonzero shift \(h_0\), the literal active atom set, both
polarizations, all reconstruction multipliers, and one inherited
global normalization \(\nu_X\) are retained.  No class may be
dropped because it is inconvenient or re-normalized after being
estimated.

\section{Why a complete tail certificate is still needed}

TPC-93 exports the complete literal low-frequency window and, under
its stated moment/truncation inequality, gives a fixed-power reserve
for the associated high-frequency reconstruction tail
\citep{WangTPC93}.  It explicitly leaves the short, resonant and
generic returned sectors to later arithmetic estimates.  TPC-102
records that closing one resonant branch still leaves the
high-frequency return, block-to-physical normalization and the
strict endpoint ledger \citep{WangTPC102}.  TPC-MVP1 packages these
requirements as Hypothesis H4 \citep{WangMVP1}.

The logical issue is a mismatch of quantifiers.  A bound for one
frequency row has the form
\[
 |\Tail_{E,\alpha}|\le B_{E,\alpha}.
\]
H4 asks for the scalar obtained only \emph{after} all outer labels
\(\alpha\), signs, masks and multipliers have been reassembled.
Summing the displayed row bounds is valid, but its total cost may
erase the saving.  Keeping the signed sum may be much better, but
then its cancellation must be proved.  The definitions below make
the two routes explicit.

\section{The literal atom partition}

For each \(X\), let \(\Omega_X\) be the finite set of inherited
physical atoms.  An atom \(\omega\) contains every outer label,
frequency, polarization and provenance key.  Let
\[
 z_X(\omega)=c_X(\omega)A_X(\omega)\in\C
\]
be its literal scalar contribution before the one global
normalization \(\nu_X>0\).  Thus
\[
 \mathcal H_X=\nu_X\sum_{\omega\in\Omega_X}z_X(\omega).
\]

\begin{definition}[Complete complement partition]
\label{def:partition}
A family of masks
\[
 \one_{\Low_X},\quad
 \one_{\High_X},\quad
 \one_{\Ultra_X},\quad
 \one_{\Short_X},\quad
 \one_{\Boundary_X},\quad
 \one_{\Collision_X}
\]
is complete if each mask is \(0\)-\(1\), their supports are pairwise
disjoint, and their sum is \(\one_{\Omega_X}\).  Empty classes are
allowed.  Further literal classes may be added, but no atom may have
coverage zero or greater than one.
\end{definition}

\begin{proposition}[Exact original-scale disintegration]
\label{prop:disintegration}
For a complete partition \(\mathcal C_X\),
\[
 \boxed{\quad
 \mathcal H_X=\sum_{C\in\mathcal C_X}\mathcal H_{X,C},
 \qquad
 \mathcal H_{X,C}
 =\nu_X\sum_{\omega\in C}z_X(\omega).
 \quad}
\]
If the masks overlap with multiplicity \(d_X(\omega)\), the displayed
sum instead equals
\(\mathcal H_X+\nu_X\sum_\omega(d_X(\omega)-1)z_X(\omega)\).
If they miss atoms, the corresponding omitted sum remains as a
literal residual.
\end{proposition}

\begin{proof}
Insert \(\sum_C\one_C(\omega)=1\) into the finite sum.  The overlap
and missing formulas follow by retaining the actual coverage
multiplicity.
\end{proof}

\begin{remark}[Classification is not estimation]
\Cref{prop:disintegration} is an exact \(\Lzero\) identity.  It does
not say that any class is small.  In particular, ``ultra-long'' and
``short'' are geometric names, not analytic bounds.
\end{remark}

\section{An optimal frequency-tail test}

Let \(G\) be a finite abelian frequency group and use a unitary
Fourier transform.  A literal row has the form
\[
 \mathcal R_X
 =\nu_X\sum_{r\in G}\widehat m_X(r)
       \overline{\widehat A_X(r)}.
\]
The multiplier \(\widehat m_X\) includes the signed reconstruction
filter.  The response \(\widehat A_X\) includes the complete
coefficient and provenance aggregation for that row.

\begin{theorem}[Sharp Hilbert tail certificate]
\label{thm:hilbert}
For every frequency set \(E\subseteq G\),
\[
 \left|\Tail_E\right|
 :=
 \nu_X\left|
 \sum_{r\in E}\widehat m_X(r)
       \overline{\widehat A_X(r)}
 \right|
 \le
 \nu_X M_E^{1/2}A_E^{1/2},
\]
where
\[
 M_E=\sum_{r\in E}|\widehat m_X(r)|^2,
 \qquad
 A_E=\sum_{r\in E}|\widehat A_X(r)|^2.
\]
The constant \(1\) is optimal.  If
\(\widehat m_X\one_E\ne0\), equality is attained by taking
\(\widehat A_X\one_E=t\widehat m_X\one_E\) for any \(t\ge0\).
\end{theorem}

\begin{proof}
This is Cauchy--Schwarz in \(\ell^2(E)\).  The stated proportional
choice makes the two vectors collinear and gives equality.
\end{proof}

\begin{corollary}[Callable high-frequency return]
\label{cor:callable}
A proved pair of complete outer estimates
\[
 M_{\High_X}\le X^{\mu+o(1)},\qquad
 A_{\High_X}\le X^{\alpha+o(1)}
\]
returns
\[
 |\Tail_{\High_X}|
 \le\nu_X X^{(\mu+\alpha)/2+o(1)}.
\]
This Hilbert certificate yields \(o(X)\) exactly when its displayed
right-hand side, with the literal \(\nu_X\), is \(o(X)\).  The actual
tail can still be smaller through cancellation not used by the
certificate.  A rowwise version is callable only after the sum of its
rowwise right-hand sides is shown to have that same property, or after
a coherent outer Hilbert estimate is proved.
\end{corollary}

\begin{remark}[What the theorem does not manufacture]
Parseval may identify \(A_E\) with a physical energy, but it does not
make that energy small.  Nor does the label ``high frequency'' force
\(M_E\) to decay.  The two quantitative inputs in
\cref{cor:callable} must come from the literal filter and arithmetic
response.
\end{remark}

\section{Sharp phase-blind obstruction}

The alternative to an energy estimate is often to bound every atom
or block separately.  The next elementary identity is the exact
limit of that strategy.

\begin{theorem}[Worst coherent completion]
\label{thm:phaseblind}
Let \(a_\omega\ge0\) be prescribed magnitudes on a finite set
\(\Omega\).  Then
\[
 \boxed{\quad
 \sup_{|\varepsilon_\omega|=1}
 \left|\sum_{\omega\in\Omega}
       \varepsilon_\omega a_\omega\right|
 =\sum_{\omega\in\Omega}a_\omega .
 \quad}
\]
Consequently, from magnitude data
\(|z_X(\omega)|\le a_\omega\) alone, the best uniform original-scale
bound is
\[
 |\mathcal H_X|
 \le\nu_X\sum_\omega a_\omega,
\]
and the constant one is sharp.
\end{theorem}

\begin{proof}
The triangle inequality gives the upper bound.  Taking every
\(\varepsilon_\omega=1\) attains it.
\end{proof}

\begin{corollary}[No support-only tail saving]
\label{cor:no-support}
Knowing only that a complement is sparse, decomposes into short
fibers, or has bounded overlap cannot give an \(o(X)\) bound unless
those facts imply
\[
 \nu_X\sum_{\omega\in\mathrm{complement}}|z_X(\omega)|=o(X).
\]
Otherwise a coefficient-sensitive signed theorem, an energy theorem
with quantitative input, or a different exact representation is
necessary.
\end{corollary}

\begin{remark}[A sharp local obstruction is not a global lower bound]
\Cref{thm:phaseblind} says that the \emph{information supplied} cannot
exclude alignment.  It does not assert that the actual arithmetic
coefficients align, and therefore proves no lower bound for the
physical tail.  Its conclusion is a stop certificate for a
phase-blind proof route, not for the prime-pair problem.
\end{remark}

\section{Complete conditional closure}

Split the complement classes into soft classes \(\Soft_X\) and
retained packet classes \(\Ret_X\).  ``Soft'' means that an
original-scale absolute or coherent estimate is actually available;
it is not a geometric synonym.

\begin{theorem}[Original-scale complement closure and H4 specialization]
\label{thm:closure}
Assume:
\begin{enumerate}[label=(\roman*),leftmargin=2.2em]
\item the masks form a complete partition in the sense of
  \cref{def:partition};
\item the complete signed low window is retained without changing
  its normalization;
\item the soft complements satisfy
  \[
  \sum_{C\in\Soft_X}|\mathcal H_{X,C}|=o(X);
  \]
\item after all reconstruction, provenance and outer reassembly
  costs, the retained complement satisfies the \emph{aggregate}
  estimate
  \[
  \left|\sum_{C\in\Ret_X}\mathcal H_{X,C}\right|
  \le X^{1-\sigma+o(1)}
  \quad\text{for some fixed }\sigma>0.
  \]
\end{enumerate}
Then the complete complement is \(o(X)\), and
\[
 \mathcal H_X=\mathcal H_{X,\Low_X}+o(X)
\]
at the original physical normalization.  A high-frequency class may
meet item (iii) or (iv) through \cref{thm:hilbert}, but only with its
complete outer energies.  In addition, item (iv) supplies the
retained-packet clause of MVP1 Hypothesis H4 only when one may take
\[
 \sigma\ge \sigma_*=\frac1{400}.
\]
An arbitrary fixed \(\sigma>0\) proves the displayed \(o(X)\)
complement closure, but does not by itself preserve the raw H4 packet
reserve.  Any later physical reassembly loss is charged separately in
H9.
\end{theorem}

\begin{proof}
Apply \cref{prop:disintegration}, then the hypotheses to the two
disjoint complement sums.  Since a fixed power saving is \(o(X)\),
their total is \(o(X)\).
\end{proof}

\begin{remark}[Per-packet estimates are insufficient by themselves]
If there are \(N_X\) retained packets and only
\(|\mathcal H_{X,C}|\le X^{1-\sigma+o(1)}\) is known separately,
triangle summation costs \(N_X\).  The theorem requires the resulting
aggregate exponent, including this cost, not the best packet
exponent quoted in isolation.
\end{remark}

\section{Finite exact regression}

The standard-library script
\texttt{experiments/tpc116\_tail\_certificate.py} uses rational
arithmetic to verify:
\begin{enumerate}[leftmargin=2em]
\item exact one-fold coverage and exact atom disintegration;
\item the explicit overlap-defect identity;
\item Cauchy--Schwarz for rational tail vectors and equality on
  proportional vectors;
\item the exact phase-blind supremum for all real sign assignments
  in finite examples; and
\item a soft/retained closure ledger that refuses incomplete masks
  and unknown class bounds.
\end{enumerate}
The committed JSON is an algebra regression only.  It contains no
computed prime-pair evidence and no asymptotic high-tail estimate.

\section{Verdict and next gate}

\begin{center}
\small
\begin{tabularx}{\textwidth}{@{}p{0.18\textwidth}p{0.39\textwidth}X@{}}
\toprule
Status & Required literal certificate & Consequence\\
\midrule
Go for \(o(X)\) closure &
Complete masks and complete outer soft/retained bounds satisfying
\cref{thm:closure} &
The complement is \(o(X)\) at the original normalization; pass the
same atoms to the TPC-18 cover audit.\\
Go for the H4 packet interface &
The preceding certificate with retained exponent
\(\sigma\ge\sigma_*=1/400\) &
The raw H4 reserve is preserved; H9 reassembly costs remain
separate.\\
Stop phase-blind route &
Only support/magnitude data are known and their total absolute mass
is not \(o(X)\) &
\Cref{thm:phaseblind} is sharp; retain signs or prove energy
cancellation.\\
Stop declaration &
An atom is missing, duplicated, or returned at a changed
normalization &
The claimed complete complement is not the original coefficient.\\
Current TPC state &
Local high-tail reserves exist, but no archived complete outer
certificate for every H4 class is supplied here &
H4 remains open; no zero loss may be entered in the endpoint ledger.\\
\bottomrule
\end{tabularx}
\end{center}

\paragraph{Established.}
The partition identity, optimal Hilbert bound, phase-blind supremum,
conditional complement closure and H4 specialization are exact
\(\Lzero\) statements.  Their
attachment to the literal TPC labels and normalization is an
\(\Lone\) interface.

\paragraph{Not established.}
We do not prove the complete outer energy or absolute mass of any
growing TPC high/ultra-tail carrier.  We do not verify every H4 atom
from a machine-readable coefficient archive.  No new \(\Ltwo\)
fixed-\(h_0\) saving, parity breaking, Hardy--Littlewood asymptotic,
prime-pair lower bound or twin-prime theorem is obtained.  The next
gate is a literal TPC-18 block cover and physical residual audit.

\begingroup
\small
\setlength{\bibsep}{2pt plus 0.3ex}
\sloppy
\bibliographystyle{plainnat}
\bibliography{references}
\endgroup
\end{document}
